\section{The Book and Its Title Page}

The small book that carries this doctrine into most English-speaking hands announces itself, on the title page of the printing currently in trade, as \emph{Trustful Surrender to Divine Providence: The Secret of Peace and Happiness}, ``By Fr.~Jean Baptiste Saint-Jure, S.J.\ and St.~Claude de la Colombière, S.J., Translated by Prof.~Paul Garvin.'' The copyright page records a nihil obstat by Rev.~John P. Sullivan and an imprimatur of Francis Cardinal Spellman, Archbishop of New York, dated 8 November 1961; original publication in 1961 by Alba House, Staten Island, under the title \emph{The Secret of Peace and Happiness}; reissue by St.~Raphael Editions of Sherbrooke, Quebec, in 1978 and 1980 and, in 1983, in cooperation with TAN Books; copyright 1980 in St.~Raphael Editions; and a 2012 TAN printing at Charlotte, North Carolina.\endnote{Front matter of the TAN Books printing of 2012, read 2026-07-25 in the publisher's own posted preview file (\texttt{tanbooks.com}, \emph{Preview\_TrustfulSurrender.pdf}), which reproduces the title leaf, copyright page, publisher's note, table of contents, and the opening pages of the text. The bibliographic and ecclesiastical facts recited here are facts, not expression. The translation itself remains in copyright and is not quoted in this article except for the two short phrases of the publisher's note discussed below, which are the object of the analysis.}

That is a great deal of publishing history for a book of about 130 pages, and it is honest as far as it goes. What it does not do---what no title page could do---is tell the reader which sentences were written by which of the two men, in what genre, or by whose hand they were selected, arranged, and joined. A short unsigned note facing the table of contents supplies the publisher's answer: the following pages were written by Saint-Jure (1588--1657) and are ``an extract from his great work entitled \emph{The Knowledge and Love of Our Lord Jesus Christ}''; to them ``has been added a further extract from the writings of St.~Claude de la Colombière (1641--1682).''\endnote{TAN printing of 2012, unnumbered leaf (p.~v), read 2026-07-25 in the publisher's preview file. The same note asserts that Saint-Jure's work ``was the constant companion during life of the saintly Curé of Ars''; that claim is reported here as the publisher's, and no independent witness for it was checked. The two quoted phrases are reproduced because the accuracy of the publisher's own description is the question under examination.} The table of contents then divides the book: Part~I, three chapters, ``by Father Jean Baptiste Saint-Jure''; Part~II, two chapters, ``by Saint Claude de la Colombière.''

This section tests that division against the sources. The test is possible because both underlying corpora are available in public-domain nineteenth-century printings, and because the English arrangement preserves enough structure to be matched against them. The result is that the publisher's note is substantially right about Part~I, substantially right about the first chapter of Part~II, and---on the evidence assembled here---wrong about the last chapter, which is not Colombière's at all.

\subsection{Part I: Saint-Jure's Chapter on Conformity}

Jean Baptiste Saint-Jure's \emph{De la connaissance et de l'amour du Fils de Dieu Notre-Seigneur Jésus-Christ} is a very large seventeenth-century work of spiritual theology, reprinted repeatedly in France into the nineteenth century; the copy used here is the Paris edition of Périsse frères, 1877, whose second volume carries the end of Book~II and Book~III through chapter~X.\endnote{Saint-Jure, \emph{De la connaissance et de l'amour du fils de Dieu notre seigneur Jésus Christ}, vol.~2 (Paris: Périsse frères, 1877); public-domain scan and OCR text consulted 2026-07-25 at the Internet Archive, item \texttt{delaconnaissance02sain} (publisher, place, and date confirmed from the item's catalogue record; the volume's running heads confirm the Book and chapter divisions cited). The date of the work's first publication in the 1630s was not verified at a witness for this article and is therefore not asserted; the 1877 printing is a re-edition, and no collation of it against any seventeenth-century printing was performed.} Book~III, chapter~VIII bears the title \emph{L'amour de Notre-Seigneur unit notre volonté à la sienne}---the love of Our Lord unites our will to his---and it opens like this:

\begin{quote}
\footnotesize Je commencerai le grand exercice de la conformité de notre volonté avec celle de Dieu Notre-Seigneur par la doctrine de l'oracle de la théologie, saint Thomas (1~p., q.~19, a.~4). Parlant de la volonté de Dieu, il dit qu'elle est la cause de toutes les choses créées.\endnote{Saint-Jure, \emph{De la connaissance}, vol.~2 (Périsse, 1877), Book~III, ch.~VIII, § I, p.~206; wording read 2026-07-25 in the Internet Archive OCR of the item cited above. Working English gloss by this article: ``I shall begin the great exercise of the conformity of our will with that of God Our Lord from the teaching of the oracle of theology, St.~Thomas (Part~I, q.~19, a.~4). Speaking of the will of God, he says that it is the cause of all created things.'' Obvious OCR artefacts have been silently normalised only where the printed reading is unambiguous.}
\end{quote}

The first chapter of the English Part~I begins from the same authority at the same locus, opening with St.~Thomas on the will of God as the cause of all that exists and carrying a footnote citing \emph{Summa}, Part~I, q.~19, a.~4.\endnote{TAN printing of 2012, Part~I, ch.~I, p.~3, with its note 1 (``St.~Thomas, Sum.\ p.~1, q.~19, a.~4; St.~Augustine, De Gen.''), read 2026-07-25 in the publisher's preview file. The correspondence of authority, locus, and argumentative function is reported; no sentence of the English translation is reproduced.} The correspondence continues through the chapter. Saint-Jure's section~VII asks \emph{en quoi nous devons pratiquer cette conformité}---in what matters we are to practise this conformity---and answers: in the natural things outside us (heat, cold, rain, hail, storms, plague, famine), in what touches us more closely, in defects of natural endowment, in sickness; section~VIII continues into death and into the degrees of grace and glory.\endnote{Saint-Jure, \emph{De la connaissance}, vol.~2, Book~III, ch.~VIII, § VII (p.~256, with its analytic head-note: ``I.~Dans les choses naturelles qui sont hors de nous.---II.~Dans celles qui nous touchent de plus près.---III.~Dans les défauts des perfections de nature.---IV.~Dans les maladies'') and § VIII (``I.~Dans la mort.---II.~Dans les vertus et les degrés de la grâce et de la gloire''); read 2026-07-25 at the witness cited above.} The English Part~I, chapter~III, subdivides the practice of conformity into fourteen numbered heads that run over the same ground in the same order: natural incidents of daily life, public calamities, family life, reverses of fortune, poverty, adversity and disgrace, defects of nature, sickness and infirmity, death and the manner of it, loss of spiritual consolation, the consequences of our sins, interior trials, spiritual favours, and a summary.\endnote{TAN printing of 2012, table of contents, Part~I, ch.~III, heads 1--14, read 2026-07-25 in the preview file. The English subdivision is finer than the French sectioning, which is what one expects of a translator working an extract into a stand-alone book; the sequence and subject matter, however, track Saint-Jure's §§ VII--VIII closely.} Part~I is thus what the publisher says it is: an extract, more finely subdivided than its original, from a single chapter of a single work.

\subsection{Part II, Chapter IV: Colombière's Sermons}

Claude de la Colombière is a different kind of author. He left no large systematic treatise on providence. What he left, and what was printed after his death, is preaching: sermons delivered largely in England during his time as chaplain to the Duchess of York, together with retreat notes and \emph{Réflexions chrétiennes}. The nineteenth-century collected \emph{Œuvres} of 1832 print the sermons in thematic order, each headed by its thesis. The fifth volume contains, in this order, sermons \emph{Sur la soumission à la volonté de Dieu} (p.~65), \emph{Sur la confiance en Dieu} (p.~85), \emph{Sur la prière} (p.~102), and, later in the volume, \emph{Sur les adversités} (p.~218).\endnote{\emph{Œuvres du R.~P. Claude de la Colombière, de la Compagnie de Jésus} (1832), vol.~5, table of sermons and the sermon texts themselves; public-domain OCR consulted 2026-07-25 at the Internet Archive, item \texttt{uvresdurpclauded05laco}. The volume's table gives each sermon with its thesis, from which the theses quoted in the body are taken verbatim.}

The theses printed with those sermons are worth setting beside the English chapter's subdivisions. \emph{Sur les adversités} is summarised in the 1832 table as: \emph{Les adversités nous sont utiles si nous sommes justes, elles nous sont nécessaires si nous sommes pécheurs}---adversities are useful to us if we are just, and necessary to us if we are sinners. The English Part~II, chapter~IV, contains a numbered head reading ``Adversity is useful for the just and necessary for sinners.'' \emph{Sur la prière} is summarised as: \emph{Nous obtenons peu par nos prières, parce que nous demandons trop peu, parce que le peu que nous demandons nous ne le demandons pas assez}---we obtain little by our prayers because we ask too little, and because the little we ask we do not ask enough. The English chapter's prayer heads run: ``Recourse to prayer,'' ``To obtain what we want,'' ``To be delivered from evil,'' ``We do not ask enough,'' ``Perseverance in prayer,'' ``Obstinate trust.''\endnote{Sermon theses from \emph{Œuvres} (1832), vol.~5, as cited; English chapter heads from the TAN table of contents, Part~II, ch.~IV, read 2026-07-25. The sermon \emph{Sur la confiance en Dieu} is preached on \emph{Fides tua te salvum fecit} (Luke~17) with the thesis that God has bound himself to help those who put their trust in him, and that nothing is more apt to engage him to do so than that trust; the sermon \emph{Sur la soumission à la volonté de Dieu} closes with an act of acceptance beginning \emph{J'accepte cette calamité, et en elle-même, et dans toutes ses circonstances}. Both were read at the same witness on the same date.} The correspondence is close enough, thesis by thesis, that Part~II's first chapter is best described as a compilation from several Colombière sermons rather than an ``extract'' from one text---which is a materially different editorial act, since a compiler choosing among sermons preached to different congregations on different subjects is making a book that its author never wrote.

\subsection{The Chapter That Changed Authors}

The last chapter of Part~II, printed under Colombière's name, is headed ``Exercise of Conformity to Divine Providence'' and divided into three parts: an act of faith, hope and charity; an act of filial submission to providence; and the usefulness of this exercise.\endnote{TAN printing of 2012, table of contents, Part~II, ch.~V, heads 1--3 (pp.~122--127), read 2026-07-25.}

Saint-Jure's chapter on conformity closes with a section headed \emph{Exercice particulier envers la providence divine}, analysed in its own head-note as: \emph{I.~Actes de foi.---II.~D'espérance.---III.~D'amour et d'estime.---IV.~D'abandon absolu}. It opens by saying that the exercise is of marvellous importance and inestimable profit; it then directs the reader to make first a great act of faith in the truth of providence, that God has a continual and most particular care of everything and of you especially---of your body, your soul, your occupations, your dwellings, your contentments, your reputation, your necessities, your health, your sicknesses, your life, your death, down to the least of your hairs, which neither falls nor even stirs without his disposition; then an act of hope for the same objects; then an act of charity, loving providence tenderly as a child its good mother; and finally, after these acts have been made and often repeated with firmness and vigour, the soul is to abandon itself totally to divine providence, to rest and as it were fall gently asleep in its arms as a child in its mother's, taking for its motto the words of David, \emph{In pace in idipsum dormiam et requiescam}, and singing further \emph{Dominus regit me, et nihil mihi deerit}.\endnote{Saint-Jure, \emph{De la connaissance}, vol.~2, Book~III, ch.~VIII, § X, pp.~280--282, read 2026-07-25 at the Internet Archive witness cited above; the summary in the body is a close working paraphrase of the French, with the Latin of Ps.~4:9 and Ps.~22:1 (Vulgate numbering) as Saint-Jure quotes them.}

The English chapter reproduces this section, in this order, with these images and these two psalms, collapsing Saint-Jure's four movements into three headings by joining faith, hope, and charity under one head and giving the fourth---absolute abandonment---the title ``Act of filial submission to providence.'' Its opening sentence states, as the French does, that the exercise is of great importance because of the advantages it confers on those who undertake it devoutly.\endnote{Correspondence established 2026-07-25 by reading the chapter in an unauthorised full-text posting of the English translation found on the open web. That posting is not an authorised witness, no wording from it is reproduced in this article, and it is recorded here only because candour requires stating how the correspondence was checked. The structural facts on which the argument rests---three heads in the stated order, the child-in-the-mother's-arms image, and the two psalms---are independently visible in the publisher's own table of contents and in Saint-Jure's French.} On the evidence, the chapter that the book attributes to Colombière is a translation of Saint-Jure.

Three cautions belong with that conclusion. First, seventeenth-century spiritual writers borrowed freely and often silently; it is not impossible that Colombière reproduced Saint-Jure's exercise somewhere in his own papers, and the present research did not exhaust his corpus. What can be said is that the text is Saint-Jure's, that it stands in Saint-Jure's chapter as its designed conclusion, and that no corresponding text was found in the Colombière volumes examined. Second, the misattribution is at the level of the compiler's parts, not of the doctrine: the exercise says nothing that the rest of the book does not say. Third---and this is the point that matters for a reader---the compiler's arrangement is invisible to the person holding the book, and its effect is to lend a canonised saint's name to pages he did not write, and to leave the actual author of much of the book, a man never raised to the altars, at the level of a supporting credit.

\subsection{Drift in the Attribution}

The drift did not stop at the compiler. Colombière was beatified in 1929 and canonised in 1992, so that printings before and after the canonisation differ on the title page between ``Blessed'' and ``Saint''---an ordinary and correct change.\endnote{Beatification 1929 (Pius~XI) and canonisation 31 May 1992 (John Paul~II) are reported at the level of standard reference works, checked 2026-07-25; no act of the Holy See was consulted directly for this article, and nothing in the argument depends on the exact dates. Older English printings of the book, and current second-hand listings of them, carry ``Blessed Claude de la Colombière'' where the 2012 TAN printing carries ``St.''} Less ordinary is what happens outside the publisher's control. Extracts of the book circulate on devotional sites and in parish handouts under the heading ``Trustful Surrender to Divine Providence, by St.~Claude de la Colombière,'' with Saint-Jure's name dropped; at least one large Catholic retailer lists the volume that way.\endnote{Observed 2026-07-25 in open-web copies and retail listings, including a widely mirrored extract PDF headed ``by St.~Claude de la Columbière'' that in fact prints both the Colombière compilation and the Saint-Jure exercise. The observation is offered as evidence of reception drift, not as a charge against any particular publisher; no listing is treated as an authority for anything.} A book of composite authorship, whose parts were selected by an unnamed hand, translated in 1961, and reissued under three imprints in four decades, has thus in ordinary use become the work of the more famous of its two named authors---who did not write the chapter most often reprinted from it.

None of this touches whether the doctrine is true. It touches what kind of witness the book is. It is not a treatise arguing a thesis; it is a devotional anthology, at two removes from both authors, addressed to readers assumed to be already instructed in the faith. Sentences in it that sound like technical theology---and there are several, of which the most important is treated later in these pages---are seventeenth-century preaching prose, and must be read as such, with the qualification the preacher himself supplies elsewhere and with the doctrinal grammar the Church supplies always.

\begin{fourstable}{Part as printed}{Attributed to}{Source identified here}{Basis and limits of the identification}
Part I, chs.~I--III & Saint-Jure & \emph{De la connaissance et de l'amour du Fils de Dieu}, Book~III, ch.~VIII, §§ I--VIII (Périsse ed., 1877, vol.~2) & Opening authority and locus (ST~I, q.~19, a.~4) identical; the practice-of-conformity heads track §§ VII--VIII in order and subject. Extract, resubdivided; not collated line by line, and no seventeenth-century printing consulted.\\
Part II, ch.~IV & Colombière & Compilation from sermons \emph{Sur la soumission à la volonté de Dieu}, \emph{Sur la confiance en Dieu}, \emph{Sur la prière}, and \emph{Sur les adversités} (\emph{Œuvres}, 1832, vol.~5) & Chapter heads reproduce the printed theses of those sermons, including ``useful for the just, necessary for sinners'' and ``we do not ask enough.'' Sermon-by-sermon collation was not performed; other Colombière material may also be present.\\
Part II, ch.~V & Colombière & Saint-Jure, same chapter, § X (\emph{Exercice particulier envers la providence divine}) & Four movements (faith, hope, love, absolute abandonment) rendered as three heads in the same order, with the same maternal image and the same two psalms. Contrary evidence would be a parallel in Colombière's own corpus, which was searched but not exhausted.\\
Whole book & Saint-Jure and Colombière & An anonymous compiler's anthology, translated by Paul Garvin (1961) & The selection, joining, and division into parts are editorial acts by an unnamed hand; the imprimatur of 1961 attests doctrinal review of the English book, not authorship or textual criticism.\\
\end{fourstable}
